\documentclass[12pt]{article}
\usepackage[T1]{fontenc}
\usepackage[utf8]{inputenc}
\usepackage{lmodern}
\usepackage{textcomp}
\usepackage{enumitem}
\usepackage{geometry}
\geometry{margin=1in}
\begin{document}
\begin{center}
Answer Key - Astronomy - Graduate Level Assessment 20 \
\small (Answers are ordered exactly as in the question paper.)
\end{center}

\section*{Multiple Choice}
\begin{enumerate}[label=\arabic*.]
\item \textbf{Answer:} C
\\ \textit{Options:} A) ancient heavily cratered highlands; B) dark lavas inside volcanic calderas; C) dark lavas filling older impact basins; D) the bright regions on the Moon
\\ \textit{Note:} The lunar maria are formed from basaltic lava flows that erupted after the formation of large impact basins, resulting in their characteristic dark appearance compared to the lighter highland regions of the Moon.
\item \textbf{Answer:} B
\\ \textit{Options:} A) To measure interstellar mass.; B) To measure galactic distances.; C) To measure galactic energy-density.; D) To measure interstellar density.
\\ \textit{Note:} Cepheid stars are relevant for astronomers because their predictable luminosity variations allow for accurate distance measurements to galaxies, acting as "standard candles" in the cosmos.
\item \textbf{Answer:} D
\\ \textit{Options:} A) the Sun mainly emits blue light.; B) the atmosphere absorbs mostly blue light.; C) molecules scatter red light more effectively than blue light.; D) molecules scatter blue light more effectively than red light.
\\ \textit{Note:} The correct answer is supported by the Rayleigh scattering phenomenon, which demonstrates that shorter wavelengths of light, such as blue, are scattered more efficiently by atmospheric molecules than longer wavelengths like red, resulting in the blue appearance of the sky.
\item \textbf{Answer:} C
\\ \textit{Options:} A) the entire planets are made mostly of metal.; B) metals condensed first in the solar nebula and the rocks then accreted around them.; C) metals sank to the center during a time when the interiors were molten throughout.; D) radioactivity created metals in the core from the decay of uranium.
\\ \textit{Note:} The correct answer is based on the principle of differentiation, where denser materials, such as metals, gravitate towards the center of a molten planetary body, leading to metallic cores in terrestrial planets.
\item \textbf{Answer:} D
\\ \textit{Options:} A) Perseus; B) Cygnus; C) Scorpius; D) Leo
\\ \textit{Note:} Leo is a zodiac constellation that is located away from the plane of the Milky Way, while the other options typically lie along or near this galactic structure.
\end{enumerate}

\section*{True / False}
\begin{enumerate}[label=\arabic*.]
\setcounter{enumi}{5}
\item \textbf{Answer:} True
\\ \textit{Options:} A) True; B) False
\\ \textit{Note:} The statement is true because the relationship between an object's actual brightness (luminosity) and its apparent brightness as observed from Earth is governed by the inverse square law, which allows astronomers to calculate the distance to the object using these two measurements.
\item \textbf{Answer:} True
\\ \textit{Options:} A) True; B) False
\\ \textit{Note:} The statement is true because Mercury's elliptical orbit causes it to be approximately 1.5 times closer to the Sun at perihelion than at aphelion, resulting in a solar radiation flux ratio of 9:4 due to the inverse square law of radiation intensity.
\item \textbf{Answer:} False
\\ \textit{Options:} A) True; B) False
\\ \textit{Note:} Differentiation is primarily driven by processes such as gravitational separation of materials based on density and thermal evolution, rather than spontaneous emission from radioactive atoms.
\item \textbf{Answer:} True
\\ \textit{Options:} A) True; B) False
\\ \textit{Note:} The statement is true because the visible part of the electromagnetic spectrum encompasses wavelengths that fall between approximately 380 nanometers (violet light) and 740 nanometers (red light), which corresponds to the range of light detectable by the human eye.
\item \textbf{Answer:} True
\\ \textit{Options:} A) True; B) False
\\ \textit{Note:} The period of rotation is the same for every point on a solid disk because all points on the disk complete one full rotation around the axis of rotation in the same amount of time, regardless of their distance from that axis.
\end{enumerate}

\section*{Long Form Response}
\begin{enumerate}[label=\arabic*.]
\setcounter{enumi}{10}
\item \textbf{Answer:} The Laniakea supercluster, which includes the Milky Way, represents a significant cosmic structure that reshapes our understanding of galactic evolution and the large-scale organization of the universe. By encompassing approximately 100,000 galaxies and spanning over 520 million light-years, Laniakea highlights the interconnectedness of galaxy clusters and their gravitational influences, challenging the traditional view of isolated galaxies. This supercluster's discovery has profound implications for the understanding of cosmic flows and the dynamics of galaxy formation, suggesting that the evolution of the Milky Way is intricately linked to its cosmic neighborhood. Furthermore, studying Laniakea allows astronomers to refine their models of dark matter and the expansion of the universe, thereby enhancing our comprehension of the fundamental processes that govern cosmic evolution.
\\ \textit{Note:} correct because it emphasizes the Laniakea supercluster's role in illustrating the interconnectivity of galaxies and clusters, thereby enhancing our understanding of cosmic evolution and challenging previous notions of isolated galactic structures.
\item \textbf{Answer:} A parsec is defined as the distance at which one astronomical unit (AU) subtends an angle of one arcsecond as viewed from Earth. This unit of measurement is significant in astronomy because it provides a practical way to express vast cosmic distances, facilitating the comparison of distances between stars and galaxies. The relationship between parsecs, AU, and angular measurements stems from the principles of parallax, where the apparent shift in position of a nearby star against a distant background is used to calculate its distance. As a result, one parsec is equivalent to approximately 3.26 light- years, emphasizing its utility in the cosmological context where distances can be exceedingly large. Understanding parsecs is essential for astronomers, as it allows for accurate distance measurements crucial for studying the structure and dynamics of the universe.
\\ \textit{Note:} correct because it accurately defines a parsec in terms of its geometric relationship to astronomical units and angular measurements, emphasizing its significance in measuring vast cosmic distances through the principle of parallax.
\end{enumerate}

\vfill
\noindent\textit{End of Answer Key}
\end{document}